% ========================================================
% ALGORITMO CONCURRENTE Y MODELO EN PROMELA
% ========================================================
\section{Algoritmo concurrente y modelo en Promela}\label{sec:promela}

\subsection{Planteamiento del algoritmo}

El algoritmo es Lloyd: asignar cada viaje al centroide más cercano y recalcular cada centroide como
la media de los viajes asignados, hasta que los centroides dejen de moverse o se agoten las
iteraciones. La versión secuencial y la concurrente comparten el mismo contrato numérico, fijado
antes de escribir código para que el benchmark compare dos implementaciones del mismo algoritmo y
no dos algoritmos distintos:

\begin{itemize}
  \item Distancia euclídea al cuadrado sobre las seis features; en caso de empate gana el centroide
        de menor índice;
  \item Inicialización k-means++ con semilla fija, ejecutada una sola vez; los centroides iniciales se
        guardan en disco con su sha256 y ambas versiones parten de ese archivo;
  \item Driterio de parada $\max_j \lVert \mu_j^{(t)} - \mu_j^{(t-1)} \rVert_\infty \le
        10^{-9} + 10^{-9} \cdot \max_j \lVert \mu_j^{(t-1)} \rVert_\infty$, o el máximo de
        iteraciones, lo que ocurra primero; el resultado registra cuál de los dos paró la corrida;
  \item Un cluster vacío conserva su centroide anterior.
\end{itemize}

La versión concurrente (Algoritmo~\ref{alg:lloyd-chunks}) parte los $N$ viajes en $C$ bloques de
tamaño $c$, los reparte por un canal entre $P$ goroutines que viven toda la corrida, y sincroniza una
sola vez por iteración, en la barrera. Cada bloque tiene sus propios acumuladores, y la reducción los
suma en orden de índice de bloque.

\begin{algorithm}[H]
\caption{Lloyd concurrente con worker pool, bloques por canal y acumuladores por bloque}
\label{alg:lloyd-chunks}
\begin{algorithmic}[1]
\Require viajes $X = \{x_1, \dots, x_N\} \subset \mathbb{R}^6$, centroides iniciales
         $\mu_1, \dots, \mu_k$, $P$ workers, tamaño de bloque $c$
\Ensure centroides finales, asignaciones $a(i)$ e inercia $J$
\State $C \gets \lceil N / c \rceil$; crear un canal \textit{trabajos} con capacidad $C$
\State Lanzar $P$ goroutines; cada una repite: recibir un bloque $b$ del canal, procesarlo, avisar
       a la barrera \Comment{el pool no se relanza por iteración}
\For{$t = 1, \dots, \textit{maxIter}$}
  \State barrera.Add($C$); enviar $b = 0, \dots, C-1$ por \textit{trabajos}
  \Statex \hspace{\algorithmicindent} \textit{cada worker, al recibir $b$:} \Comment{$\mu$ es de solo lectura}
  \State \hspace{\algorithmicindent} $s_{bj} \gets \mathbf{0}$, $n_{bj} \gets 0$, $J_b \gets 0$
  \ForAll{$x_i$ en el bloque $b$}
    \State $j^{*} \gets \arg\min_j \lVert x_i - \mu_j \rVert^2$; \ $a(i) \gets j^{*}$
    \State $s_{bj^{*}} \gets s_{bj^{*}} + x_i$; \ $n_{bj^{*}} \gets n_{bj^{*}} + 1$;
           \ $J_b \gets J_b + \lVert x_i - \mu_{j^{*}} \rVert^2$
  \EndFor
  \State \hspace{\algorithmicindent} barrera.Done()
  \State barrera.Wait() \Comment{\texttt{sync.WaitGroup}: recién ahora se puede escribir $\mu$}
  \State $s_j \gets \sum_{b=0}^{C-1} s_{bj}$, \ $n_j \gets \sum_{b=0}^{C-1} n_{bj}$,
         \ $J \gets \sum_b J_b$ \Comment{en orden de $b$, siempre el mismo}
  \State $\mu_j \gets s_j / n_j$ para todo $j$ con $n_j > 0$
  \If{$\max_j \lVert \mu_j^{(t)} - \mu_j^{(t-1)} \rVert_\infty$ cumple la tolerancia} \textbf{salir} \EndIf
\EndFor
\State Una pasada más de asignación con los $\mu$ finales, también en paralelo, para $a(i)$ y $J$
\State Cerrar \textit{trabajos}; esperar a que las $P$ goroutines terminen
\State \Return $\mu$, $a$, $J$
\end{algorithmic}
\end{algorithm}

\subsection{Qué se modela}

El modelo en Promela (Listado~\ref{lst:pml}) no es el K-means: no tiene distancias, centroides
reales ni punto flotante. Es el esqueleto de sincronización del Algoritmo~\ref{alg:lloyd-chunks},
con un dominio deliberadamente pequeño, cuatro puntos, dos bloques, dos workers y dos iteraciones,
para que Spin pueda recorrer \emph{todos} los entrelazados posibles. Lo que importa está completo:
el canal que reparte bloques entre workers persistentes, los acumuladores privados por bloque, la
barrera al final de cada iteración (la variable \texttt{pendientes} es el \texttt{WaitGroup}) y la
publicación de centroides nuevos, que solo puede ocurrir cuando no queda ningún lector.

Se modelan dos iteraciones a propósito. Los errores de reutilizar mal la barrera o de no reiniciar
los acumuladores no aparecen en la primera iteración: aparecen en la segunda.

\subsection{Correspondencia entre la implementación y el modelo}

La abstracción conserva las decisiones de concurrencia y elimina solamente el cálculo numérico.
La Tabla~\ref{tab:mapeo-promela} hace explícita esa correspondencia. Esta trazabilidad es importante:
una propiedad demostrada sobre variables que no representan fielmente el programa no permite sacar
conclusiones sobre su diseño.

\begin{table}[H]
\caption{Correspondencia entre Go y el modelo Promela}\label{tab:mapeo-promela}
\small
\begin{tabularx}{\textwidth}{@{}L{4.1cm}L{4.2cm}Y@{}}
\toprule
\textbf{Implementación en Go} & \textbf{Abstracción en Promela} & \textbf{Aspecto conservado} \\
\midrule
\texttt{trabajos chan int} & \texttt{chan trabajos = [C] of \{byte\}} &
  Reparto dinámico de índices de bloque entre workers persistentes \\
\addlinespace
Goroutines del pool & Dos procesos \texttt{worker()} &
  Consumo concurrente y no determinista de los bloques \\
\addlinespace
\texttt{parciales[c]} & \texttt{conteo[c]} &
  Propiedad de cada acumulador por un único bloque durante la ronda \\
\addlinespace
\texttt{barrera.Add}, \texttt{Done} y \texttt{Wait} & \texttt{pendientes = C}, decrementos y guarda
  \texttt{pendientes == 0} & La reducción empieza solo después de terminar todos los bloques \\
\addlinespace
Lectura de \texttt{cent} y llamada a \texttt{actualizar} &
  \texttt{lectores} y \texttt{escribiendo} & Exclusión temporal entre asignación y publicación \\
\addlinespace
Rangos disjuntos de observaciones & \texttt{seen[i]} &
  Cobertura individual, sin omisiones ni duplicados \\
\bottomrule
\end{tabularx}
\end{table}

No se modelan los valores de los centroides ni las sumas de seis dimensiones porque las propiedades
estudiadas no dependen de ellos. \texttt{conteo[c]} representa el efecto observable de procesar un
bloque; \texttt{total} representa la reducción de sus parciales. La versión en Go admite un último
bloque más corto, mientras que el dominio acotado usa $N=4$ y $C=2$, de modo que ambos bloques tienen
tamaño \texttt{TAM=2}. Esta simplificación reduce el espacio de estados, pero no introduce una prueba
sobre los límites del último bloque.

\subsection{Ejecución de una ronda en Promela}

\texttt{init} crea los dos workers dentro de un bloque \texttt{atomic}; esto evita entrelazados
irrelevantes durante la creación, no durante el procesamiento. Al comenzar una ronda, reinicia
\texttt{seen}, \texttt{total} y, en el mutante, \texttt{compartido}; luego fija
\texttt{pendientes=C} antes de enviar los dos índices al canal. El orden es el mismo que
\texttt{WaitGroup.Add} antes de encolar en Go.

La selección \texttt{trabajos ? idx} queda habilitada cuando existe un mensaje. SPIN decide qué
worker lo recibe y explora las alternativas de planificación. Cada worker incrementa
\texttt{lectores}, reinicia el acumulador del bloque, marca sus dos puntos y finalmente decrementa
\texttt{pendientes}. En Promela, una asignación simple es un paso indivisible; por eso las variables
de instrumentación cuentan eventos sin necesitar otro cerrojo.

La expresión aislada \texttt{pendientes == 0} no es una comprobación que pueda fallar: es una
\emph{guarda}. Si el contador todavía es positivo, \texttt{init} queda bloqueado y solo los workers
pueden avanzar. Cuando llega a cero, el coordinador comprueba que no queden lectores, activa
\texttt{escribiendo}, reduce los parciales y valida las invariantes de la ronda. Después repite el
ciclo sin recrear el pool, igual que la implementación.

\subsection{Qué se verifica}

\begin{table}[H]
\caption{Propiedades verificadas por Spin sobre el modelo}\label{tab:propiedades}
\small
\begin{tabularx}{\textwidth}{@{}L{5.6cm}Y@{}}
\toprule
\textbf{Propiedad} & \textbf{Cómo se expresa} \\
\midrule
Cada punto se procesa exactamente una vez por iteración
  & \texttt{assert(seen[i] == 1)} para todo $i$ al final de cada iteración \\
\addlinespace
Los centroides no se escriben mientras alguien los lee
  & \texttt{assert(escribiendo == 0)} en cada worker al leer; \texttt{assert(lectores == 0)} en el
    coordinador antes de publicar \\
\addlinespace
La reducción conserva los puntos & \texttt{assert(total == N)} tras sumar los acumuladores \\
\addlinespace
Ausencia de deadlock & Estados finales válidos, que \texttt{pan} comprueba por su cuenta; los
  workers bloqueados esperando trabajo llevan la etiqueta \texttt{end}, que declara legítima esa
  espera \\
\bottomrule
\end{tabularx}
\end{table}

Sobre la primera propiedad conviene una advertencia. Sumar los conteos no alcanza: procesar el
punto A dos veces y omitir el punto B da exactamente el mismo total, así que
\texttt{assert(total == N)} pasaría con un error de cobertura. Por eso se cuenta por punto, con el
arreglo \texttt{seen}.

\subsection{El mutante: por qué el cero errores significa algo}

Un modelo que da cero errores solo prueba algo si es capaz de encontrar errores. Para comprobarlo,
el mismo archivo se compila con \texttt{-DMUTANTE}, que reemplaza el acumulador privado por uno
compartido sin exclusión mutua y separa la lectura de la escritura
(\texttt{tmp = compartido; tmp++; compartido = tmp}). La separación es necesaria: si lectura y
escritura fueran un solo enunciado, Promela lo ejecutaría de forma indivisible y la carrera quedaría
escondida. Este es el mismo incremento perdido que el curso estudia con el algoritmo de
\textcite{benari2006} en las primeras semanas, ahora dentro del problema real.

\subsection{Resultados}

Los Listados~\ref{lst:spin-ok} y~\ref{lst:spin-mut} son la salida de \texttt{pan} sobre el modelo
correcto y sobre el mutante con SPIN 6.5.2. Se usó búsqueda exhaustiva, no \emph{bitstate}; la
reducción de orden parcial que informa la salida elimina conmutaciones equivalentes sin convertir la
búsqueda en aproximada. El modelo correcto termina con \texttt{errors: 0} tras almacenar 2\,407
estados y ejecutar 3\,229 transiciones, con profundidad máxima 180. El mutante encuentra la
violación de \texttt{compartido == 4} en el paso 205 y alcanza profundidad 208 antes de detener la
búsqueda en el primer error; por eso su salida dice \texttt{Search not completed}, como corresponde
a una ejecución destinada a producir un contraejemplo, y no a validar el modelo mutante.

El rastro se reproduce paso a paso con \texttt{spin -t -p -DMUTANTE kmeans.pml}. El
\texttt{Makefile} de \texttt{tp/spin/} comprueba en cada corrida que el modelo correcto produzca
cero errores y el mutante exactamente uno, para que el modelo no deje de probar lo que dice probar
sin que nadie lo note.

\lstinputlisting[style=terminal, caption={Salida de \texttt{pan} sobre el modelo correcto},
                 label={lst:spin-ok}, lastline=22]{generado/spin-correcto.txt}

\lstinputlisting[style=terminal, caption={Salida de \texttt{pan} sobre el mutante con acumulador compartido},
                 label={lst:spin-mut}, lastline=16]{generado/spin-mutante.txt}

\subsection{Hasta dónde llega esta prueba}

La conclusión vale para este modelo y estos parámetros. No es una demostración del programa en Go
para cualquier tamaño: lo que descarta es toda una clase de errores, pérdida o duplicación de
puntos, escritura de centroides durante la lectura, barrera mal reutilizada, sobre todos los
entrelazados de ese dominio. Eso es exactamente lo que ningún test puede hacer, porque un test
observa un entrelazado por corrida. Por eso el curso pide las dos cosas y este trabajo hace las dos:
Spin verifica el diseño (esta sección) y \texttt{go test -race} verifica la implementación
(Sección~\ref{sec:go}). Las propiedades de progreso, que toda iteración termine, requieren declarar
hipótesis de \emph{fairness} y quedan para el Entregable 3.

\subsection{El modelo}

\lstinputlisting[language=Promela, firstline=32,
                 caption={\texttt{tp/spin/kmeans.pml}, sin el comentario de cabecera},
                 label={lst:pml}]{../../spin/kmeans.pml}
